Medical fine-tuning changes how vision-language models use an image, but it is unclear
whether it makes them rely less on the question text. We measure this with per-layer
attention knockout, which blocks image-to-answer or question-to-answer attention one
layer at a time. The probe is applied as a pre-specified before/after instrument to
Gemma-3-4B and its medical fine-tune MedGemma-4B on the VQA-RAD radiology benchmark
(250 samples: discovery 0000--0099, confirmation 0100--0249).

On the confirmation split, the pre-specified primary test holds: the median per-sample
question-dominance ratio falls from 5.36 to 4.15 ($p{=}0.006$, Cliff's $\delta{=}0.18$),
a small effect relative to the model's language prior. Four of five discovery candidates
replicate (one only at the raw level). They show a more decisive first answer token,
smaller closed yes/no logit margins, and lower image-to-answer information flow (IF AUC,
replicated at the raw level only). A discovery-only ``image enters later'' shift
does \emph{not} replicate and is reported as a null. Closed-ended accuracy is higher for
the medical model (77.2\% vs.\ 44.3\%, $p{=}3.9{\times}10^{-5}$; both are forced-choice
logit verdicts, not generated answers), but the base model sits
at the always-``yes'' rate, below the majority-class baseline, so part of this gap is a
baseline effect. The medical vision encoder differs sharply from the base (linear CKA
0.155). Whether the change lies in the encoder or the decoder is left open and is the
subject of follow-up work.
